% GENERATED FILE. Edit canonical lessons, then run npm run generate:books.
% canonical-source-hash: fnv1a64:f4aab3b0f45341cf
% canonical-lessons: KA-C508-medulu, KA-C508-svasakosa, KA-C508-mutrapinda, KA-C508-karulu, KA-C508-snayu

\chapter{Inside the Body}
\label{ch:ka-inside-the-body}
\hlchaptermodality{508}

\begin{chapteropening}
\textbf{By the end of this chapter:} \emph{I can say the brain, a lung, a kidney, the intestine and a muscle.}

Complete the last lesson of chapter 508: I can say the brain, a lung, a kidney, the intestine and a muscle.
\end{chapteropening}

\section[\texorpdfstring{\kn{ಮೆದುಳು} (meduḷu)}{ಮೆದುಳು (meduḷu)}]{\kn{ಮೆದುಳು} (meduḷu) — the brain}

\label{lesson:KA-C508-medulu}



\begin{quote}
Before the new one: say the Kannada for hiccups, then the Kannada for asthma.
\end{quote}

\subsection*{You'll want to know: \kn{ಮೆದುಳು}}
\textbf{\kn{ಮೆದುಳು}} — \emph{meduḷu} — \textquotedblleft{}the brain\textquotedblright{}.

\begin{quote}
\kn{ಮೆದುಳು} \kn{ದೇಹದ} \kn{ಎಲ್ಲಾ} \kn{ಕೆಲಸಗಳನ್ನು} \kn{ನಡೆಸುತ್ತದೆ}.
\end{quote}

\begin{quote}
\emph{meduḷu dēhada ellā kelasagaḷannu naḍesuttade.}
\end{quote}

\begin{quote}
The brain runs all the work of the body.
\end{quote}

Also written \kn{ಮಿದುಳು}.

\begin{cousinweb}
\textbf{\kn{ಮೆದುಳು}} (\emph{meduḷu}) is native Dravidian: DEDR 5062 lists Kannada \emph{miduḷ} and \emph{meduḷ}, \textquotedblleft{}brain, marrow\textquotedblright{}, beside Telugu \emph{medaḍu}.
\end{cousinweb}

\begin{grammarlens}[title={What you've built}]
One more word for describing what you lived and felt.
\end{grammarlens}

\subsection*{Your turn}
\begin{itemize}
\raggedright
  \item \emph{Say it:} \emph{meduḷu}
  \item \emph{Say it:} \emph{meduḷu}, once more
  \item \emph{Say it:} \emph{meduḷu dēhada ellā kelasagaḷannu naḍesuttade.} — The brain runs all the work of the body.
  \item \emph{Recall:} say the Kannada for hiccups, then the Kannada for asthma, then say \emph{meduḷu} again
\end{itemize}

\begin{tcolorbox}[breakable,colback=teal!4,colframe=teal!35!black,title={Before you move on}]
What does \kn{ಮೆದುಳು} mean? (\textquotedblleft{}The brain\textquotedblright{}.) Say it once more.
\end{tcolorbox}
\section[\texorpdfstring{\kn{ಶ್ವಾಸಕೋಶ} (śvāsakōśa)}{ಶ್ವಾಸಕೋಶ (śvāsakōśa)}]{\kn{ಶ್ವಾಸಕೋಶ} (śvāsakōśa) — a lung}

\label{lesson:KA-C508-svasakosa}



\begin{quote}
Before the new one: say the Kannada for asthma, then the Kannada for the brain.
\end{quote}

\subsection*{You'll want to know: \kn{ಶ್ವಾಸಕೋಶ}}
\textbf{\kn{ಶ್ವಾಸಕೋಶ}} — \emph{śvāsakōśa} — \textquotedblleft{}a lung\textquotedblright{}.

\begin{quote}
\kn{ಹೊಗೆಯಿಂದ} \kn{ಶ್ವಾಸಕೋಶಕ್ಕೆ} \kn{ಹಾನಿ} \kn{ಆಗುತ್ತದೆ}.
\end{quote}

\begin{quote}
\emph{hogeyinda śvāsakōśakke hāni āguttade.}
\end{quote}

\begin{quote}
Smoke harms the lungs.
\end{quote}

Plural: \kn{ಶ್ವಾಸಕೋಶಗಳು}, \emph{śvāsakōśagaḷu}.

\kn{ಶ್ವಾಸ}, breath + \kn{ಕೋಶ}, a case: the case that holds the breath.

\begin{cousinweb}
\textbf{\kn{ಶ್ವಾಸಕೋಶ}} (\emph{śvāsakōśa}) is a learned compound built from Sanskrit \emph{śvāsa}, \textquotedblleft{}breath\textquotedblright{} (CDIAL 12769), and \emph{kośa}, \textquotedblleft{}a case, a sheath\textquotedblright{} (CDIAL 3539).
\end{cousinweb}

\begin{grammarlens}[title={What you've built}]
One more word for describing what you lived and felt.
\end{grammarlens}

\subsection*{Your turn}
\begin{itemize}
\raggedright
  \item \emph{Say it:} \emph{śvāsakōśa}
  \item \emph{Say it:} \emph{śvāsakōśa}, once more
  \item \emph{Say it:} \emph{hogeyinda śvāsakōśakke hāni āguttade.} — Smoke harms the lungs.
  \item \emph{Recall:} say the Kannada for asthma, then the Kannada for the brain, then say \emph{śvāsakōśa} again
\end{itemize}

\begin{tcolorbox}[breakable,colback=teal!4,colframe=teal!35!black,title={Before you move on}]
What does \kn{ಶ್ವಾಸಕೋಶ} mean? (\textquotedblleft{}A lung\textquotedblright{}.) Say it once more.
\end{tcolorbox}
\section[\texorpdfstring{\kn{ಮೂತ್ರಪಿಂಡ} (mūtrapiṇḍa)}{ಮೂತ್ರಪಿಂಡ (mūtrapiṇḍa)}]{\kn{ಮೂತ್ರಪಿಂಡ} (mūtrapiṇḍa) — a kidney}

\label{lesson:KA-C508-mutrapinda}



\begin{quote}
Before the new one: say the Kannada for the brain, then the Kannada for a lung.
\end{quote}

\subsection*{You'll want to know: \kn{ಮೂತ್ರಪಿಂಡ}}
\textbf{\kn{ಮೂತ್ರಪಿಂಡ}} — \emph{mūtrapiṇḍa} — \textquotedblleft{}a kidney\textquotedblright{}.

\begin{quote}
\kn{ಸಾಕಷ್ಟು} \kn{ನೀರು} \kn{ಕುಡಿದರೆ} \kn{ಮೂತ್ರಪಿಂಡ} \kn{ಚೆನ್ನಾಗಿರುತ್ತದೆ}.
\end{quote}

\begin{quote}
\emph{sākaṣṭu nīru kuḍidare mūtrapiṇḍa cennāgiruttade.}
\end{quote}

\begin{quote}
If you drink enough water, the kidneys stay healthy.
\end{quote}

Doctors and patients also say \kn{ಕಿಡ್ನಿ}.

\begin{cousinweb}
\textbf{\kn{ಮೂತ್ರಪಿಂಡ}} (\emph{mūtrapiṇḍa}) is a learned compound built from Sanskrit \emph{mūtra}, \textquotedblleft{}urine\textquotedblright{} (CDIAL 10234), and \emph{piṇḍa}, \textquotedblleft{}a lump, a ball\textquotedblright{} (CDIAL 8168).
\end{cousinweb}

\begin{grammarlens}[title={What you've built}]
One more word for describing what you lived and felt.
\end{grammarlens}

\subsection*{Your turn}
\begin{itemize}
\raggedright
  \item \emph{Say it:} \emph{mūtrapiṇḍa}
  \item \emph{Say it:} \emph{mūtrapiṇḍa}, once more
  \item \emph{Say it:} \emph{sākaṣṭu nīru kuḍidare mūtrapiṇḍa cennāgiruttade.} — If you drink enough water, the kidneys stay healthy.
  \item \emph{Recall:} say the Kannada for the brain, then the Kannada for a lung, then say \emph{mūtrapiṇḍa} again
\end{itemize}

\begin{tcolorbox}[breakable,colback=teal!4,colframe=teal!35!black,title={Before you move on}]
What does \kn{ಮೂತ್ರಪಿಂಡ} mean? (\textquotedblleft{}A kidney\textquotedblright{}.) Say it once more.
\end{tcolorbox}
\section[\texorpdfstring{\kn{ಕರುಳು} (karuḷu)}{ಕರುಳು (karuḷu)}]{\kn{ಕರುಳು} (karuḷu) — the intestine; the gut}

\label{lesson:KA-C508-karulu}



\begin{quote}
Before the new one: say the Kannada for a lung, then the Kannada for a kidney.
\end{quote}

\subsection*{You'll want to know: \kn{ಕರುಳು}}
\textbf{\kn{ಕರುಳು}} — \emph{karuḷu} — \textquotedblleft{}the intestine; the gut\textquotedblright{}.

\begin{quote}
\kn{ಹಸಿವಾದಾಗ} \kn{ಕರುಳು} \kn{ಸದ್ದು} \kn{ಮಾಡುತ್ತದೆ}.
\end{quote}

\begin{quote}
\emph{hasivādāga karuḷu saddu māḍuttade.}
\end{quote}

\begin{quote}
When you are hungry, the gut rumbles.
\end{quote}

Also the seat of tender feeling: \kn{ಕರುಳಿನ} \kn{ಸಂಬಂಧ}, a blood tie.

\begin{cousinweb}
\textbf{\kn{ಕರುಳು}} (\emph{karuḷu}) is native Dravidian: DEDR 1274 lists Kannada \emph{karuḷ}, \textquotedblleft{}an entrail, the bowels; love\textquotedblright{}, beside Malayalam \emph{karaḷ}.
\end{cousinweb}

\begin{grammarlens}[title={What you've built}]
One more word for describing what you lived and felt.
\end{grammarlens}

\subsection*{Your turn}
\begin{itemize}
\raggedright
  \item \emph{Say it:} \emph{karuḷu}
  \item \emph{Say it:} \emph{karuḷu}, once more
  \item \emph{Say it:} \emph{hasivādāga karuḷu saddu māḍuttade.} — When you are hungry, the gut rumbles.
  \item \emph{Recall:} say the Kannada for a lung, then the Kannada for a kidney, then say \emph{karuḷu} again
\end{itemize}

\begin{tcolorbox}[breakable,colback=teal!4,colframe=teal!35!black,title={Before you move on}]
What does \kn{ಕರುಳು} mean? (\textquotedblleft{}The intestine; the gut\textquotedblright{}.) Say it once more.
\end{tcolorbox}
\section[\texorpdfstring{\kn{ಸ್ನಾಯು} (snāyu)}{ಸ್ನಾಯು (snāyu)}]{\kn{ಸ್ನಾಯು} (snāyu) — a muscle}

\label{lesson:KA-C508-snayu}



\begin{quote}
Before the new one: say the Kannada for a kidney, then the Kannada for the intestine.
\end{quote}

\subsection*{You'll want to know: \kn{ಸ್ನಾಯು}}
\textbf{\kn{ಸ್ನಾಯು}} — \emph{snāyu} — \textquotedblleft{}a muscle\textquotedblright{}.

\begin{quote}
\kn{ವ್ಯಾಯಾಮದಿಂದ} \kn{ಸ್ನಾಯುಗಳು} \kn{ಗಟ್ಟಿಯಾಗುತ್ತವೆ}.
\end{quote}

\begin{quote}
\emph{vyāyāmadinda snāyugaḷu gaṭṭiyāguttave.}
\end{quote}

\begin{quote}
Exercise makes the muscles strong.
\end{quote}

Plural: \kn{ಸ್ನಾಯುಗಳು}, \emph{snāyugaḷu}.

\kn{ಸ್ನಾಯು} \kn{ಸೆಳೆತ}, a muscle cramp.

\begin{cousinweb}
\textbf{\kn{ಸ್ನಾಯು}} (\emph{snāyu}) is borrowed whole from Sanskrit \emph{snāyu}, \textquotedblleft{}tendon, muscle, sinew\textquotedblright{} (CDIAL 13793).
\end{cousinweb}

\begin{grammarlens}[title={What you've built}]
One more word for describing what you lived and felt.
\end{grammarlens}

\subsection*{Your turn}
\begin{itemize}
\raggedright
  \item \emph{Say it:} \emph{snāyu}
  \item \emph{Say it:} \emph{snāyu}, once more
  \item \emph{Say it:} \emph{vyāyāmadinda snāyugaḷu gaṭṭiyāguttave.} — Exercise makes the muscles strong.
  \item \emph{Recall:} say the Kannada for a kidney, then the Kannada for the intestine, then say \emph{snāyu} again
\end{itemize}

\begin{tcolorbox}[breakable,colback=teal!4,colframe=teal!35!black,title={Before you move on}]
What does \kn{ಸ್ನಾಯು} mean? (\textquotedblleft{}A muscle\textquotedblright{}.) Say it once more.
\end{tcolorbox}
